\section{Matched Retail Experiments}
\label{sec:experiments}

\subsection{Data and evaluation design}
\label{sec:setup}\label{sec:experimental-setup}\label{sec:datasets}

We use daily sales from M5, Rohlik~v2, and Favorita. The datasets differ
substantially in the frequency of zero sales and in the available covariates
(Table~\ref{tab:t51}). M5 and Rohlik are used in full for the preliminary
baseline runs. Favorita is limited to the 30,000 highest-volume series because
the full data did not fit the available Azure instance. The main paired
comparison then selects the same 100 high-volume series from each dataset for
every model.

\begin{table}[htbp]
\centering
\caption{Datasets used in the retail experiments.}
\label{tab:t51}
\footnotesize
\begin{tabularx}{\textwidth}{lrrrlX}
\toprule
Dataset & Full series & Baseline cap & Days & Approx. zero fraction & Covariates used by LightGBM \\
\midrule
M5 & 30,490 & 30,490 & 1,941 & 70\% & Price, calendar events, weekday, month \\
Rohlik~v2 & 5,390 & 5,390 & 1,402 & 1.2\% & Price, holidays, closures, weekday, month \\
Favorita & 174,685 & 30,000 & 1,684 & 15--25\% & Promotions, oil, holidays, transactions, calendar \\
\bottomrule
\end{tabularx}
\end{table}

The last 20\% of dates form the evaluation period. Forecasts are produced from
non-overlapping rolling origins at horizons 7, 14, and 28. All models in a
dataset use identical series and origins. Our main accuracy measure is
aggregate \WAPE{}, $\sum |y-\hat y|/\sum |y|$, calculated from the stored
predictions. This avoids the instability of averaging ratios for low-volume
series. MAE and sMAPE are retained in the replication output, and M5 receives
additional scaled-error checks. We do not compare the latter with the official
M5 competition score because our panel and rolling-origin design differ from
the competition task.
\label{sec:evalprotocol}\label{sec:fm_metrics}\label{sec:fm-m5-metric}

Uncertainty is estimated by resampling series. The same bootstrap draws are
used for every model within a dataset, preserving the dependence created by
common observations and baselines. We report 95\% percentile intervals for
the log error ratio $\ell$. A negative interval indicates lower error for the
foundation model throughout the interval.

\subsection{Models and implementation}
\label{sec:modelfamilies}\label{sec:fm}\label{sec:fm_scope}
\label{sec:fm_extraction}\label{sec:fm_local}\label{sec:fm-consumer-box}

Seasonal Naive repeats the value from seven days earlier. The LightGBM models
are global Tweedie regressions using demand lags, rolling averages, calendar
variables, and the dataset-specific covariates in Table~\ref{tab:t51}. The
recursive specification feeds earlier predictions into later steps; the
direct specification fits a separate model for each step. Both are evaluated
with fixed hyperparameters and with a 10-trial search at a validation origin.
An additional direct variant scales its tree budget with the horizon. The
standard training window is the most recent 365 days. M5 also receives a
richer feature set and a separate 200-trial tuning check.

Five foundation models were feasible on the available MacBook Air M3 with
16\,GB unified memory: Chronos-Bolt-Tiny, Moirai~2.0-Small, TiRex, Chronos-2,
and TimesFM~2.5. Their parameter counts range from approximately 9 to 200
million. Each is used zero-shot, without retail covariates or parameter
tuning. TiRex runs in batches on the CPU; the other models use their supported
MPS or CPU path. Thus the comparison is deliberately pragmatic, but it is not
a study of GPU throughput, fine-tuning, or the largest available checkpoints.

The conventional models were run on a four-vCPU Azure instance with 32\,GB
RAM, whereas the foundation models were run locally. Runtime and energy
records are therefore reported within each setting rather than converted into
a single cost ranking. Code, environments, data manifests, predictions, and
run ledgers are versioned with the replication package.
\label{sec:compute}\label{sec:reproducibility}\label{sec:fm_repro}

\subsection{The LightGBM comparator}
\label{sec:results}\label{sec:crossdataset}\label{sec:perdataset}
\label{sec:baseline}\label{sec:baseline-reliability}
\label{sec:m5results}\label{sec:dirvsr}\label{sec:m5distance}
\label{sec:rohlik}\label{sec:favorita}\label{sec:synthesis}
\label{sec:threats}\label{sec:costconsequences}\label{sec:takeaways}

LightGBM improves on Seasonal Naive in the full-workload preliminary runs: by
28\% on M5 \WRMSSE{} at horizon 7, by 9\% on Rohlik \WAPE{}, and by 17\% on
Favorita \WAPE{}. The choice between direct and recursive forecasting is less
stable. Table~\ref{tab:lgbm-protocol} reports the comparison on the 100-series
panels used for the foundation models. Direct forecasting is slightly worse
at short horizons on M5, consistently worse on Rohlik, and better only at
horizon 28 on M5 and Favorita. Increasing the direct model's tree budget does
not improve these panel results.

\begin{table}[htbp]
\centering
\caption{Error of direct LightGBM relative to recursive LightGBM on the
matched panels. Positive values mean that direct forecasting has higher
aggregate \WAPE{}.}
\label{tab:lgbm-protocol}\label{tab:t52}\label{tab:t52b}\label{tab:t58}
\small
\begin{tabular}{lrrr}
\toprule
Dataset & $h=7$ & $h=14$ & $h=28$ \\
\midrule
M5 & $+0.4\%$ & $+1.1\%$ & $-2.1\%$ \\
Rohlik~v2 & $+3.9\%$ & $+4.3\%$ & $+4.0\%$ \\
Favorita & $+4.6\%$ & $+2.2\%$ & $-4.9\%$ \\
\bottomrule
\end{tabular}
\end{table}

The corresponding full-workload, per-series results are more favourable to
direct forecasting on M5. This difference arises from both the sampled series
and the aggregation of \WAPE{}. It cautions against describing either
multi-step strategy as intrinsically stronger. For the comparisons below, we
therefore select the lowest-error observed baseline separately for each
dataset and horizon. Detailed baseline grids are reported in
Appendix~\ref{app:baseline-detail}.

\subsection{Foundation-model results}
\label{sec:fmresults_pointer}\label{sec:fm_results}

Table~\ref{tab:sourceb-paired} summarises the 45 comparisons. Chronos-2 has lower
aggregate \WAPE{} than the best observed baseline in every dataset--horizon
cell, with intervals below zero. The estimated reduction ranges from about
4\% on M5 at horizons 14 and 28 to about 8\% on Favorita at horizons 7 and 14.
TiRex also has a lower point estimate in all nine cells, although intervals
include zero for M5 and Favorita at horizon 28. Moirai is lower in eight cells
and tied in the ninth. TimesFM performs well on M5 and Rohlik but is close to
the LightGBM baseline on Favorita. Chronos-Bolt-Tiny is near parity on M5 and
worse on the other two datasets.

\begin{table}[htbp]
\centering
\caption{Summary of foundation-model comparisons with the best observed
conventional baseline across three horizons. The final column counts bootstrap
intervals lying entirely below zero.}
\label{tab:sourceb-paired}
\small
\begin{tabular}{llrrr}
\toprule
Dataset & Model & Lower error & Range of $\ell$ & Intervals below zero \\
\midrule
M5 & Chronos-2 & 3/3 & $-0.062$ to $-0.038$ & 3/3 \\
M5 & Chronos-Bolt-Tiny & 0/3 & 0.006 to 0.032 & 0/3 \\
M5 & Moirai~2.0-Small & 3/3 & $-0.068$ to $-0.033$ & 3/3 \\
M5 & TimesFM~2.5 & 3/3 & $-0.061$ to $-0.030$ & 3/3 \\
M5 & TiRex & 3/3 & $-0.047$ to $-0.013$ & 2/3 \\
\addlinespace
Rohlik & Chronos-2 & 3/3 & $-0.079$ to $-0.061$ & 3/3 \\
Rohlik & Chronos-Bolt-Tiny & 0/3 & 0.008 to 0.041 & 0/3 \\
Rohlik & Moirai~2.0-Small & 3/3 & $-0.034$ to $-0.016$ & 2/3 \\
Rohlik & TimesFM~2.5 & 3/3 & $-0.060$ to $-0.040$ & 3/3 \\
Rohlik & TiRex & 3/3 & $-0.061$ to $-0.037$ & 3/3 \\
\addlinespace
Favorita & Chronos-2 & 3/3 & $-0.086$ to $-0.057$ & 3/3 \\
Favorita & Chronos-Bolt-Tiny & 0/3 & 0.162 to 0.190 & 0/3 \\
Favorita & Moirai~2.0-Small & 2/3 & $-0.024$ to 0.003 & 2/3 \\
Favorita & TimesFM~2.5 & 0/3 & 0.017 to 0.025 & 0/3 \\
Favorita & TiRex & 3/3 & $-0.049$ to $-0.025$ & 2/3 \\
\bottomrule
\end{tabular}
\end{table}

The replication files contain all 45 cell-level estimates and intervals.

Model size does not explain the ordering in this range. The 120-million
parameter Chronos-2 is consistently strong, but the smaller Moirai and TiRex
often match or improve on the 200-million parameter TimesFM. Nor is there a
single horizon pattern shared by the models. The more useful distinction is
between datasets: Favorita is the closest contest for Moirai and TimesFM,
whereas Chronos-2 retains a clearer advantage.

\subsection{Sensitivity of the Chronos-2 result}

Selecting the lowest-error LightGBM variant separately in each cell protects
against a weak comparator, but it also makes the comparison data-dependent.
We therefore repeated the Chronos-2 contrast against tuned recursive LightGBM
in all nine cells. The direction remains favourable to Chronos-2 and every
interval remains below zero (Table~\ref{tab:prespecified-tuned-recursive}).

\begin{table}[htbp]
\centering
\caption{Chronos-2 compared with the same prespecified baseline, tuned
recursive LightGBM, in every dataset--horizon cell.}
\label{tab:prespecified-tuned-recursive}
\small
% Generated by analysis/source_b_prespecified_baseline_sensitivity.py
\begin{tabular}{llrrr}
\toprule
Dataset & $h$ & FM WAPE & Baseline WAPE & $\ell$ [95\% CI] \\
\midrule
Favorita & 7 & 0.248 & 0.271 & -0.085 [-0.108, -0.064] \\
Favorita & 14 & 0.255 & 0.277 & -0.086 [-0.112, -0.061] \\
Favorita & 28 & 0.270 & 0.301 & -0.108 [-0.131, -0.086] \\
M5 & 7 & 0.294 & 0.312 & -0.062 [-0.077, -0.049] \\
M5 & 14 & 0.307 & 0.332 & -0.078 [-0.095, -0.062] \\
M5 & 28 & 0.339 & 0.368 & -0.083 [-0.102, -0.065] \\
Rohlik & 7 & 0.180 & 0.198 & -0.095 [-0.115, -0.077] \\
Rohlik & 14 & 0.185 & 0.206 & -0.109 [-0.141, -0.075] \\
Rohlik & 28 & 0.193 & 0.209 & -0.079 [-0.105, -0.054] \\
\bottomrule
\end{tabular}

\end{table}

On M5, increasing the LightGBM search from 10 to 200 trials also leaves
Chronos-2 ahead (Table~\ref{tab:m5-200trial}). The 200-trial comparison chooses
the lower-error recursive or direct variant after observing the test results,
so it is a deliberately demanding sensitivity check rather than a
pre-registered estimator. Its intervals condition on that choice.

\begin{table}[htbp]
\centering
\caption{M5 sensitivity to a 200-trial LightGBM search. Negative $\ell$
favours Chronos-2.}
\label{tab:m5-200trial}
\small
\begin{tabular}{rlrrr}
\toprule
$h$ & Selected LightGBM & Baseline \WAPE{} & Chronos-2 \WAPE{} &
$\ell$ (95\% bootstrap CI) \\
\midrule
7 & tuned recursive & 0.313 & 0.294 & $-0.063$ [$-0.077$, $-0.050$] \\
14 & tuned recursive & 0.330 & 0.307 & $-0.073$ [$-0.091$, $-0.058$] \\
28 & tuned direct & 0.360 & 0.339 & $-0.061$ [$-0.089$, $-0.032$] \\
\bottomrule
\end{tabular}
\end{table}

The main M5 panel contains high-volume items and has a median zero fraction of
0.108. A second panel, stratified by volume and zero frequency, has a median
zero fraction of 0.773. Against the 10-trial tuned LightGBM baseline,
Chronos-2's log ratios are $-0.116$, $-0.160$, and $-0.147$ at horizons 7,
14, and 28; all three bootstrap intervals exclude zero. A bottom-level,
dollar-weighted scaled-error calculation on the main panel gives the same
direction (Table~\ref{tab:m5-panel-scaled}). This measure resembles the
scaling logic of M5 but is not the official hierarchical \WRMSSE{}.

\begin{table}[htbp]
\centering
\caption{M5 scaled-error sensitivity for Chronos-2 on the matched panel.}
\label{tab:m5-panel-scaled}
\small
% Generated by analysis/source_b_m5_panel_scaled_metric.py
\begin{tabular}{rlrrr}
\toprule
$h$ & Best baseline & FM RMSSE & Baseline RMSSE & $\ell$ [95\% CI] \\
\midrule
7 & lightgbm\_cov & 0.649 & 0.682 & -0.050 [-0.063, -0.036] \\
14 & lightgbm\_cov & 0.701 & 0.735 & -0.047 [-0.061, -0.032] \\
28 & lightgbm\_tuned\_direct & 0.772 & 0.795 & -0.029 [-0.052, -0.009] \\
\bottomrule
\end{tabular}

\end{table}

Together, these checks make it unlikely that the observed Chronos-2 advantage
is caused only by the initial LightGBM specification or by the high-volume M5
sample. They do not establish performance on the full M5 hierarchy, on larger
Rohlik or Favorita panels, or against retailer-specific production systems.

\subsection{Computational scope}
\label{sec:costenvelope}

The archived local sweep totals 241.6 wall-clock hours, but its rows combine
different workloads and two hardware settings. We therefore use runtime and
energy values only to describe the measured scenarios. In particular,
TimesFM's CPU implementation is slow in the local setting, while batching is
important for TiRex. These observations should not be extrapolated to GPU
serving or very large batches. Moirai~2.0-Small also has a non-commercial
licence, which may matter more for deployment than its measured runtime.

The main remaining threats are the small 100-series panels, the use of
univariate foundation models against LightGBM models with covariates, possible
overlap between public datasets and pre-training data, and the absence of
high-budget LightGBM checks outside the main M5 panel. These issues are taken
up in Section~\ref{sec:discussion} and Appendix~\ref{app:threats}.
\label{sec:fm_threats}\label{sec:table59}
